\section{Explicit subspaces for the remaining intervals}\label{app:windows}

On $34$ of the $102$ intervals of Theorem~\ref{thm:profile-D} the corank bound comes from the counting
inequality with an explicit rational subspace. These intervals lie in $(3N,16.2N)$. The subspaces were found
numerically, as optimal constant subspaces for the bound below on each interval. They are dense rational
subspaces with no evident closed form, so they are listed in the formalization. Only the integer identities
and the interval conditions described below enter the proof.

\begin{theorem}\label{thm:profile}
Let $v_1,\dots,v_t\in\Z^s$ span a subspace $V\subset\Q^s$ of dimension $t$. For $j$ in a finite set $\mathcal J$
let $B_j$ be an integer matrix with $s$ columns and $r_j\le t$ an integer such that there are $t-r_j$ linearly
independent vectors $k\in\Z^t$ with $B_j\bigl(\sum_ak_av_a\bigr)=0$. Let $\mathcal E\ne0$ be the product of a
nonzero $t\times t$ minor of $(v_1,\dots,v_t)$ and, for each $j$, of a nonzero maximal minor of these kernel
vectors. Let $p\nmid\mathcal E$ be an odd prime with $H<p^2$, write $q=ap+r$ with $0\le r<p$, let $V_p=\F_p^t$ be
the space of coordinates of $V$ modulo $p$, and let $m$ and $w_j$ ($j\in\mathcal J$) be nonnegative integers.
Suppose that there are linear maps $\tau_{j,k}\colon V_p^m\to V_p$ ($k<w_j$) and $\iota\colon V_p^m\to\F_p^n$ such
that, on the common kernel of the maps $B_j\circ\tau_{j,k}$, the map $\iota$ is injective and satisfies the
hypothesis of Theorem~\ref{thm:count}. Then
\[
 \nu_N(p)\ge m\,t-\sum_{j\in\mathcal J}w_j\,r_j .
\]
\end{theorem}

\begin{proof}
Since $p\nmid\mathcal E$, the reductions of $v_1,\dots,v_t$ are linearly independent, so $V_p$ is the
reduction of $V$ in the coordinates of this basis. The kernel vectors of $B_j$ satisfy an integer identity,
so their reductions lie in the kernel of $B_j|_{V_p}$. They remain independent because their minor is a unit
modulo $p$. Hence $\operatorname{rank}(B_j|_{V_p})\le r_j$. The map $\Phi=(B_j\circ\tau_{j,k})_{j,k}$ therefore
has rank at most $\sum_jw_jr_j$, and its kernel has dimension at least $mt-\sum_jw_jr_j$.
Theorem~\ref{thm:count}, applied to $\ker\Phi$ with no further conditions, gives the bound.
\end{proof}

In the applications $m$ and the $w_j$ are affine functions of $N$ and $p$ on an interval of primes, while $V$
and the $B_j$ are fixed. So the bound is linear on the interval. The row combinations are given by
decompositions modulo $p$ (Definition~\ref{def:slot}) whose data $x_d$ range over $V_p$. The maps $\tau_{j,k}$
are Taylor coefficients at $X=-1,0,1$ and coefficients at $\infty$. Admissibility of the data is required only
on the common kernel of these conditions.

\emph{The integer identities.} For each of the $34$ intervals the formalization lists integer vectors
$v_1,\dots,v_t$ spanning $V$, a set of $t$ pivot columns with the adjugate of the corresponding minor, and for
each map $B_j$ its integer matrix, the bound $r_j$, $t-r_j$ kernel vectors in the coordinates of $V$, and again
a pivot set with an adjugate. A single Boolean function of these data, evaluated by the Lean kernel, checks
the shapes, the identities $\mathrm{adj}\cdot v_{\mathrm{piv}}=\delta\cdot I$ and the products
$B_j\bigl(\sum_ak_av_a\bigr)=0$ in integers. Its soundness theorem gives the hypotheses of
Theorem~\ref{thm:profile} on $V$ and the $B_j$, with $\mathcal E$ the product of the minors. No prime enters this
check.

\emph{The interval conditions.} The hypotheses on the conditions and on $\iota$ depend on $N$ and $p$. An
interval is described by finitely many linear inequalities in $N$ and $(p-1)/2$, and on each of finitely many
pieces of the interval the numbers $m$ and $w_j$ are affine in the same variables. On each piece a verified
procedure checks four things. The orders of the Taylor conditions are nonnegative and fit the interval. The
prime exceeds $59$. The map $\iota$ is injective on the common kernel and lands in the admissible data. The
values at the moment columns are given there by a polynomial in the column of degree less than $q\bmod p$,
so that Theorem~\ref{thm:count} applies. Each required inequality is derived from the inequalities of the
interval by nonnegative combinations and integer rounding, that is, by Chv\'atal--Gomory cuts. The soundness
theorem of the procedure gives the hypotheses of Theorem~\ref{thm:profile} at every odd prime $p>59$ of the
interval with $p\nmid\mathcal E$. For large $N$ all primes of the interval are of this kind. For $33$ of the
$34$ intervals the hypotheses are established in this way, and for the remaining one by a direct proof of
the same statement. Since the intervals lie in $(3N,16.2N)$, where $p>q$, we have $a=0$ in
Theorem~\ref{thm:count}.

\section{The bounds for the energy method}\label{app:field}

This appendix gives the trial sequences of Section~\ref{sec:energy}, the certified bounds that they satisfy, and
the bound for $W_0$.

\emph{Trial sequences.} For the constrained case we use a trial sequence with $40$ Chebyshev coefficients and a
finite number of geometric terms, with the multiplier $\lambda_1=1.68493357$. For the case with doubled
interaction we use trial sequences of the form \eqref{eq:trial}, with $\Kcoef$ Chebyshev coefficients, which are
dyadic rationals, and $\Ngeo$ geometric terms. Of these, $\Ncomplex$ have complex $z_j$ and complex $r_j$. We
place the complex charges at the $41$ equally spaced points $x_0=-0.16+0.008j$ $(0\le j\le40)$, all with
$\epsilon_0=10^{-2}$. The real terms sit at points $\pm z$ with $z$ increasing from $0.8$ to $0.999997$. The
numbers $z_j$ are rationals with at most twelve decimals. The norms of the geometric terms have closed forms:
two terms $(z_j,r_j)$ and $(z_l,r_l)$ contribute
$-\frac12\operatorname{Re}\bigl(r_jr_l\log(1-z_jz_l)+r_j\bar r_l\log(1-z_j\bar z_l)\bigr)$ to $\|u\|^2$.

\emph{Covers by cells.} We must bound three functions of one variable from above on their whole domains: the
field $X$ on $[-1,1]$, the field $Y$ on $[0,1]$, and the field $F_1$ on $g_*\le|x|\le1$. With $x=\cos\theta$,
each is a constant plus a cosine series with $\Kcoef$ terms plus finitely many terms
$\operatorname{Re}\,w\log(1-ae^{im\theta})$. We cover each domain by cells $[x_c-h',x_c+h']$ whose endpoints are
multiples of $2^{-40}$. On each cell we bound the function by its Taylor expansion at $x_c$. The cells are
grouped in blocks. For each block the derivatives of the polynomial part are computed exactly at the centre
of the block up to order $70$, and the remainder beyond order $70$ is bounded by a single constant for the
block, multiplied by $h'^{70}/70!$. The Taylor expansion at the centre of a cell is obtained from that of the
block by an exact shift. The logarithmic terms are expanded to order $5$, with remainders bounded from the
exact minimum of $|1-ae^{im\theta}|$ on the cell. A few cells at $x=0$ and $x=\pm1$ treat the singular terms
separately. The covers of $X$, $Y$ and $F_1$ have \Nblocks{} blocks and \Ncells{} cells. Each cell is
checked by exact rational and integer arithmetic.

\emph{The bounds.} The covers give
\[
 \sup X\le\SupX,\quad \sup Y\le\SupY,\quad \sup_{g_*\le|x|\le1}F_1\le\ResCap .
\]
So $M_X+M_Y=\FieldCap$, and the norms satisfy $2\|u\|^2+\frac12\|v\|^2\le0.7737$. The first coefficient satisfies
$u_1\ge9/50$, which is used near $x=0$. For the points of the second kind a one-variable estimate gives
$\sup_{0<y<1}F_2(y)\le-77/100$ for $u_1\ge9/50$. For the constrained case, with the weights
$(\alpha_A,\beta_h,\gamma_H-\frac12,\lambda_1,2g_0/n_0)=(11/29,3/29,49/58,\lambda_1,5/29)$, the bounds
$\sup X_1\le-1.15854$ and $\sup Y\le-1.32981$ are obtained by isolating the zeros of the numerators of the
derivatives, which are polynomials with integer coefficients. The norm bound is $n_1=0.67761$. With $a_0=41/58$,
$a_1=76/29$ and $U=\Utarget$ these give the constants stated after Lemma~\ref{lem:energy-ineq}. At the rate $U_0$
the inequality of the constrained case still holds with the slope $c_2$, and that of the case with doubled
interaction holds for $\alpha\ge\frac1{16}$ with the slope $c_2$.

\emph{The bound for $W_0$.} For a fixed radius the function $y\mapsto\max_{\Gamma_r}\Phi_y$ is a maximum of
functions affine in $y$, hence convex. The entropy $h_2$ is concave. On a cell $[y_1,y_2]\subset[\frac12,1]$ we
replace $h_2$ by its tangent at the midpoint. The sum is then convex, and its supremum on the cell is
attained at an endpoint. On $[0,\frac12]$ the term $\tilde h_2=\log2$ is constant. We cover $[0,1]$ by $14$
cells, each with a rational radius $r$. At each endpoint the maximum of $\Phi_y$ on the circle and on the
segments is bounded by a cover of the contour by small arcs and intervals, in exact arithmetic. This proves
\[
 \sup_{0\le y\le1}\bigl(w(y)+58\tilde h_2(y)\bigr)\le W_0=\WZero .
\]
The supremum is attained near $y=0.62$, with radius about $0.6$.

\section{The trial measures of Section~\ref{sec:signed}}\label{app:trials}

A trial measure for the concave majorant of \S\ref{ssec:reflection} consists of a reference measure for the
free nodes with a piecewise polynomial density, trial sequences for the rows, the columns and the numerator
of Theorem~\ref{thm:heine}, and the measure $\mathsf m$ of Lemma~\ref{lem:LGlow}. The trial sequences have $120$
Chebyshev coefficients and geometric tails with complex poles, as in \eqref{eq:trial}. We use three trial
measures, at $\alpha=0$, $\frac1{200}$ and $\frac1{16}$. The constant of Theorem~\ref{thm:near1} interpolates
between them in the same way as the fields.

At each of the three trial measures we bound the suprema of six fields: the fields of the negative and the
positive free nodes, the column field, and the fields of the prescribed negative nodes, the prescribed
positive nodes and the nodes of $\mathcal S$ in \eqref{eq:expandS}. Each field is a polynomial plus finitely many
terms $P(x)\log|x-c|$ and $\operatorname{Re}\,\varrho\log(x-z)$ with complex $z$ and $\varrho$, plus a constant.

The suprema are verified on covers by cells, as in Appendix~\ref{app:field}. The endpoints of the domains are
treated by Laurent expansions and logarithmic end intervals. The polynomial part of a field has degree about
$120$. On a block of cells we shift its Taylor expansion exactly to a fixed order and bound the remaining
coefficients by one table shared by the block, which keeps the integers small. The smallest margin of the
inequality of \S\ref{ssec:signedcert} is $0.037$ per unit of $N^2$.
